\documentclass[../../main.tex]{subfiles}
\graphicspath{{../../image/}{../../figure/}}
\begin{document}

\section{真非线性系统的上确界界}\label{section:ch13-上确界界}

由取 $\tau=0$ 的\eqref{align:ch13-osc-local}，立即可得 $U_h$ 的 $L^\infty$ 模界：
\begin{align}
\sup_{(-\infty,\infty)}|U_h(\cdot,t)|\leqslant\sup_{(-\infty,\infty)}|U_0(\cdot)|+c_5TV_{(-\infty,\infty)}U_0(\cdot).
\label{align:ch13-sup-naive}
\end{align}
该估计不如所断言的\eqref{align:ch13-h-sup}强，因为除上确界外还涉及初始数据的总变差。即便如此，把\eqref{align:ch13-sup-naive}与\cref{section:ch13-总变差界}中建立的估计\eqref{align:ch13-h-tv}、\eqref{align:ch13-h-lip}结合，仍可援引\cref{section:ch13-紧性与相容性}的结果，从而推出初值问题\eqref{align:ch13-ivp}解 $U$ 的存在性，它是一列近似解的极限；见\eqref{align:ch13-convergence}。显然，由\eqref{align:ch13-h-tv}、\eqref{align:ch13-h-lip}，$U$ 满足\eqref{align:ch13-tv-bound}与\eqref{align:ch13-lip-bound}。至此我们已验证了\cref{theorem:ch13-小总变差存在性}的全部断言，除\eqref{align:ch13-sup-bound}外。该估计作为稳定性陈述颇具意义。此外，在能处理大总变差初始数据、但仍需小振幅以求解 Riemann 问题并应用波相互作用估计的特殊情形中，它发挥着有用作用。因而建立给出\eqref{align:ch13-sup-bound}的估计\eqref{align:ch13-h-sup}十分重要。

首先注意，对配备 Riemann 不变量坐标系的系统，\eqref{align:ch13-h-sup}是\cref{theorem:ch13-sup指数界}的直接推论，因而\eqref{align:ch13-tv-small}中的 $\delta_1$ 无需很小，只要\eqref{align:ch13-product-small}中的 $\delta_2$ 很小即可。此情形下证明之所以简单，是因为相互作用估计\eqref{align:ch13-int-main}中缺少二次项，即 $c_3=0$。相反，在缺乏这种特殊结构的系统中，二次阶相互作用项使情形复杂化。\eqref{align:ch13-h-sup}的证明取决于二次项的特殊形式，如\cref{chapter:ch09-容许波扇与riemann问题}所示，
它涉及 $DF$ 特征向量的 Lie 括号。该分析过于繁复，此处不拟完整重现，仅给出主要思想的梗概。读者可在\cref{section:ch13-注记}引用的文献中找到细节。

证明的一般策略受\cref{section:ch13-glimm泛函}中思想的启发，该节最终导向\cref{theorem:ch13-glimm递减}与\cref{theorem:ch13-sup指数界}。将与网格曲线 $I$ 相联系两个泛函 $\mathscr{R}$ 与 $\mathscr{P}$，其中 $\mathscr{R}(I)$ 度量 $U_h$ 在 $I$ 上的振幅，而 $\mathscr{P}(I)$ 估计未来波相互作用可能如何影响该振幅。为精确度量振幅，必须计及同一特征族激波与稀疏波的相互抵消。因而我们须对振幅求和，而非对波强度求和。相应地，对任意（有限个）波序列，比如说 $M$ 个波、振幅 $\xi=(\xi_1,\cdots,\xi_M)$，我们联系数
\begin{align}
|\xi|=\sum_{j=1}^n\left|\sum_{j\text{--波}}\xi_L\right|,
\label{align:ch13-amplitude-tally}
\end{align}
其中第二求和遍历 $L=1,\cdots,M$ 中该序列的第 $L$ 个波为 $j$ 波的那些指标。然后定义
\begin{align}
\mathscr{R}(I)=\sup|\xi|,
\label{align:ch13-R-def}
\end{align}
其中上确界取遍穿越 $I$ 的所有连续波序列，``连续''指序列中占据相邻位置的任意两个波之间夹着 $U_h$ 的一个常状态。稍加思考即可见，只要 $\mathscr{L}(I)$ 充分小，$\mathscr{R}(I)$ 就度量 $U_h$ 在 $I$ 上的振幅。

从 $I$ 推进到其后继时，$\mathscr{R}$ 的值因两个原因而变化：首先，波以不同速度传播，波序列发生穿越与重排（但注意，同一特征族波的相对顺序必然保持）；其次，波振幅因波相互作用而变化，如\cref{chapter:ch09-容许波扇与riemann问题}所示。
可以证明，波强度中三阶及更高阶相互作用的影响可如\cref{theorem:ch13-sup指数界}证明中那样粗略估计。然而，波强度中二次阶相互作用的影响更为显著，须更精确地估计。这可通过把\cref{chapter:ch09-容许波扇与riemann问题}中的二次项实现为新的\emph{虚拟波}（virtual waves）
而与实际波一并计及，从而有效完成。

上述泛函 $\mathscr{P}$ 有助于估计未来波相互作用的影响，它由下述步骤构造。对任意穿越网格曲线 $I$ 的连续波序列，联系一族波序列，视其为该序列的``后代''。后代波序列由其``亲代''序列通过以下两种运算导出：(a) 亲代序列中波的允许重排，例如亲代序列中占据第 $K$ 位的 $k$ 波与占据第 $L$ 位的 $\ell$ 波，若 $k>\ell$ 且 $K<L$，则交换位置；(b) 插入由亲代序列中波相互作用可能产生的任意虚拟波。后代序列的精确构造涉及重大技术工程，此处不拟进行，可在\cref{section:ch13-注记}引用的文献中找到。对任意网格曲线 $I$，令
\begin{align}
\mathscr{P}(I)=\sup|\xi|,
\label{align:ch13-P-def}
\end{align}
其中上确界现取遍穿越 $I$ 的所有连续波序列的后代族之并。

只要总变差很小，$\mathscr{P}$ 实际上与 $\mathscr{R}$ 等价：
\begin{align}
\mathscr{R}(I)\leqslant\mathscr{P}(I)\leqslant[1+c_9\mathscr{L}(I)]\mathscr{R}(I).
\label{align:ch13-P-equiv}
\end{align}
\eqref{align:ch13-P-equiv}证明的思想如下。回忆 $\mathscr{L}(I)$ 与 $\mathscr{R}(I)$ 的主要区别在于：前者 统计 穿越波的（正）强度，而后者对穿越波的（带符号）振幅求和，从而允许同一特征族中异号波（即激波与稀疏波）之间抵消。考察振幅为 $\zeta$ 的单个 $j$ 波与若干 $k$ 波的相互作用。由于同一特征族中的波保持相对顺序，$k$ 波与该 $j$ 波的相互作用将连续出现，从而产生的虚拟波也以相同顺序出现。进而，每当 $k$ 波的振幅符号交替时，相应的 Lie 括号项亦然。因此，虚拟波与其亲代波经历相同的抵消，从而该 $j$ 波与 $k$ 波相互作用对 $\mathscr{P}(I)$ 的贡献为 $O(1)|\zeta|\mathscr{S}(I)$ 阶。此类相互作用对 $\mathscr{P}(I)$ 的总贡献为 $O(1)\mathscr{L}(I)\mathscr{S}(I)$ 阶，由此得\eqref{align:ch13-P-equiv}。详细证明相当冗长，可见于文献。

下一步是证明，若 $J$ 是如\cref{figure:ch13-网格曲线}所示的网格曲线 $I$ 的直接后继，则
\begin{align}
\mathscr{P}(J)\leqslant\mathscr{P}(I)+c_{10}\mathscr{R}(I)\mathscr{D}(\Delta_s^r).
\label{align:ch13-P-change}
\end{align}
证明思想如下。穿越 $J$ 的波序列是穿越 $I$ 的波序列的重排，其中进入 $\Delta_s^r$ 的波在通过``西南''边与``东南''边进入后，离开 $\Delta_s^r$ 时交换相对位置。进而，从 $I$ 推进到 $J$ 时，由进入 $\Delta_s^r$ 的波相互作用产生的虚拟波被转化为实际波，体现于离开 $\Delta_s^r$ 的波中。详细证明同样冗长，应查阅文献。

由\eqref{align:ch13-P-equiv}，可在不违背不等式的前提下，于\eqref{align:ch13-P-change}右端以 $\mathscr{P}(I)$ 代替 $\mathscr{R}(I)$。因此，迭代论证可得，若 $J$ 是 $I$ 的任意后继，则
\begin{align}
\mathscr{P}(J)\leqslant\prod[1+c_{10}\mathscr{D}(\Delta_s^r)]\,\mathscr{P}(I),
\label{align:ch13-P-product}
\end{align}
其中乘积遍历夹在曲线 $I$ 与 $J$ 之间的所有菱形 $\Delta_s^r$。

现假设 $4\kappa\mathscr{L}(I)\leqslant1$，并诉诸\cref{theorem:ch13-glimm递减}。把\eqref{align:ch13-P-product}、\eqref{align:ch13-D-sum}与\eqref{align:ch13-P-equiv}结合，得到
\begin{align}
\mathscr{R}(J)\leqslant\exp[c_9\mathscr{L}(I)+c_{10}\mathscr{L}(I)^2]\,\mathscr{R}(I),
\label{align:ch13-R-exp}
\end{align}
由此即得所需的\eqref{align:ch13-h-sup}。

\end{document}